\label{ferramenta}

Este Capítulo descreve a metodologia utilizada no desenvolvimento de uma ferramenta para automatizar a geração de histórias de usuários a partir de um conjunto de requisitos éticos, desde as fases iniciais, como a obtenção do conjunto de dados, passando pela preparação dos dados, até à fase final de modelagem e avaliação da qualidade do modelo. Para identificar as características da ferramenta, foram utilizados os resultados da revisão sistemática de literatura do Capítulo \ref{rsl}.

\section{Motivação da Ferramenta}
A revisão sistemática da literatura, detalhada no Capítulo \ref{rsl}, apresentou um conjunto de características que embasaram a proposta da ferramenta desenvolvida neste trabalho. Em primeiro lugar, é importante reconhecer que os princípios éticos tendem a ser constantes \cite{rees-muller, khan, balasubramaniam}, e ainda assim as implementações existentes são lentas devido à falta de padronização e às abordagens não específicas, o que torna difícil incorporar os princípios ao processo de desenvolvimento \cite{khan, kemell, agbese, correa}. Outro problema identificado é o desenvolvimento monofocal de ferramentas. Em alguns casos, as ferramentas disponíveis nem sequer abordam o contexto completo do único princípio em questão. Essas descobertas demonstram que, atualmente, é necessário implementar uma infraestrutura que consiste em um conjunto de soluções para o desenvolvimento ético da IA durante todo o ciclo de vida. Isso introduz mais um obstáculo à operacionalização dos requisitos éticos. 

Para contornar os problemas identificados, a ferramenta proposta deve atender a alguns requisitos essenciais. Primeiramente, para superar estes desafios, é fundamental que a ferramenta proposta seja projetada para se integrar intuitivamente aos fluxos de trabalho de times de desenvolvimento já existentes, permitindo uma adoção sem dificuldades significativas. Além disso, é fundamental que ela seja baseada em padrões claros e consistentes, facilitando a incorporação dos princípios éticos de maneira uniforme e menos suscetível a variações, superando a falta de padronização que torna as implementações lentas e complexas. Além disso, a ferramenta deve tentar integrar o conjunto de princípios éticos de maneira abrangente e profunda, evitando uma abordagem monofocal e rasa. Por fim, é importante que, no caso da criação de uma infraestrutura de desenvolvimento ético de \acrshort{IA}, a ferramenta seja facilmente acoplada.

Kelleher \& Tierney \cite{kelleher} constataram que desenvolvedores de \acrshort{IA} passam 80\% do seu tempo total nas fases iniciais do projeto. Portanto, é possível deduzir que uma ferramenta que atenda aos requisitos mencionados acima e seja aplicada nos estágios iniciais do projeto representa a opção mais adequada. Dessa forma, a ferramenta proposta opera traduzindo requisitos éticos de IA em histórias de usuários éticas, facilitando assim a implementação prática de princípios éticos em sistemas de IA e preenchendo as lacunas entre a teoria e a prática. 
A ferramenta permite que as equipes de desenvolvimento traduzam os princípios éticos em tarefas transparentes, permitindo assim que os princípios sejam abordados durante todo o ciclo de vida do sistema e que sejam discutidos durante os estágios iniciais, onde o trabalho é mais concentrado. 
Além disso, ela aborda o desenvolvimento monofocal de princípios éticos, permitindo que os vários princípios do conjunto definido sejam abordados por meio de sua integração com as histórias de usuários. Com estas características definidas, a ferramenta pode então ser formalizada.
 
\section{Visão Geral}
A ferramenta é formalizada como um sistema baseado em LLM que converte requisitos éticos de \acrshort{IA} de alto nível em histórias de usuários éticas, com o objetivo de auxiliar equipes de desenvolvimento na incorporação prática de princípios éticos em seus sistemas de \acrshort{IA}. A ferramenta utiliza as técnicas de \textit{retrieval-augumented generation} e \textit{fine-tuning} para gerar predições. Ao adicionar RAGs ao sistema, a ferramenta combina o uso de um modelo de linguagem com o mecanismo de recuperação de informações em bancos de dados e documentos específicos \cite{rag, rag2}. Isso permite a busca de referências em padrões éticos, diretrizes legais e casos de uso relevantes para aumento de precisão e contextualização das gerações do modelo. O diagrama da ferramenta pode ser observado na Figura \ref{fig:diag}

% GERAR UM DIAGRAMA DA FERRAMENTA

% APRESENTAR TEMPLATE DO PEKKA E EXPLICAR A ESCOLHA

% DEFNIR AS ENTRADAS E SAÍDAS DA FERRAMENTA

\section{\emph{Dataset}}
A obtenção dos dados no contexto de desenvolvimento de soluções éticas em \acrshort{IA} assume uma complexidade adicional se comparada aos softwares convencionais, pois os dados são menos acessíveis. O dataset utilizado se encontra neste Trabalho\footnote{\url{https://zenodo.org/records/10938484}} e foi construído por meio da adição de duas bases de dados: \textit{PROMISE} \cite{CanedoM20} e \textit{Passenger Flow} \cite{pass}. A base de dados \textit{PROMISE} \cite{CanedoM20} é composta por requisitos funcionais e requisitos não funcionais para sistemas de software convencionais, não relacionados a \acrshort{IA}. Em comparação, a base de dados \textit{Passenger Flow} \cite{pass} é constituída por requisitos éticos em \acrshort{IA}, obtidos a partir de um estudo acadêmico aplicado num contexto industrial. Este dataset foi compartilhado para com os autores e no momento ainda não pode ser referenciado.

É importante notar que a coluna que contém os requisitos apresentados no formato de história de usuário não está incluída no dataset. Por este motivo, esta coluna foi acrescentada ao dataset para o treinamento do modelo. Portanto, dado o tempo e a dificuldade de obtenção dos dados necessários, foi decidido utilizar um processo de \textit{data augumentation} por meio de LLMs. De acordo com Ding et al. \cite{data-aug}, \textit{data augumentation} envolve a adoção de métodos destinados a  reforçar a eficácia do modelo através do aumento dos dados de treino por meio de dados gerados sinteticamente, evitando assim a necessidade de esforços adicionais de coleta de dados. Ao relacionar este processo com a utilização de LLMs, Ding et al. \cite{data-aug} explicam que:
\blockquote{Do ponto de vista dos dados, a utilização de LLMs para aumentar os dados representa uma estratégia viável para ultrapassar as limitações relacionadas a coleta de dados, facilitando a criação de conjuntos de dados sintéticos de alta qualidade que, em certos casos, podem revelar-se de maior valor do que os dados selecionados por humanos.}

Assim, foi realizada a transcrição manual de um subconjunto menor, contendo dez requisitos, que foram reescritos como histórias de usuário e posteriormente avaliados por pares. Após esta avaliação, os restantes requisitos foram submetidos ao processo definido de \textit{data augumentation} por meio de LLMs. Este processo utilizou como entrada a definição de cada princípio ético e as dez histórias de usuário geradas manualmente, garantindo a relevância das histórias de usuário resultantes para com o tópico.

\section{Preparação dos Dados}

O processo de pré-processamento de dados envolve a limpeza e a preparação do texto para etapas subsequentes do fluxo de \emph{machine learning}. Os tipos de pré-processamento vão de acordo com o objetivo final do sistema, não sendo uma fórmula universal \cite{haddi}. Neste estudo, os algoritmos de pré-processamento foram desenvolvidos utilizando a linguagem de programação Python e estão disponíveis no repositório do Github\footnote{\href{https://github.com/joaorossi15/requirements-to-us}{https://github.com/joaorossi15/requirements-to-us}}. 

Com o objetivo de preparar os dados, foram exploradas as técnicas descritas na literatura e foi identificado um conjunto de técnicas principais de pré-processamento \cite{pre-processing, angiani}. Dentro deste conjunto de técnicas, foram selecionadas três principais para utilização neste modelo. A primeira técnica é composta da transformação dos caracteres maiúsculos em caracteres minúsculos com o objetivo de normalizar as entradas \cite{angiani}, e é apresentado no \emph{Listing} \ref{f:min}. A segunda técnica é a remoção de espaços redundantes para garantir a formatação correta \cite{angiani} e o código relacionado a ela pode ser observado no \emph{Listing} \ref{f:rem}. De acordo com Palomino \& Aider \cite{pre-processing}, "a tokenização de um fluxo de caracteres em uma sequência de elementos semelhantes a palavras é um dos componentes mais críticos do pré-processamento de texto". Portanto, as entradas foram transformadas em tokens, como se pode observar no \emph{Listing} \ref{f:tokenization}.

\begin{itemize}
    \item Transformação do conjunto de caracteres da entrada em caracteres minúsculos: Substituição de caracteres maiúsculos por minúsculos para normalizar a formatação dos textos de entrada.
    \begin{code}
    \captionof{listing}{Função de Transformação em Minúsculo}
    \label{f:min}
    \begin{mintedbox}{python}
    def text_to_lower(text: str) -> str:
        return text.lower()
    \end{mintedbox}
    \end{code}

    \item Remoção de espaços redundantes: Remoção de espaços desnecessários para normalizar a formatação dos textos de entrada.
    \begin{code}
    \captionof{listing}{Função de Remoção de Espaços Redundantes}
    \label{f:rem}
    \begin{mintedbox}{python}
    def pad_multiple_spaces(text: str) -> str:
        return re.sub( r"\s+"," ",text)
    \end{mintedbox}
    \end{code}

    \item Tokenização dos caracteres de entrada: O processo de tokenização é de extrema importância para o processamento de linguagem natural. Modelos de LLM utilizam tokens que representam unidades semânticas básicas para compreender melhor a semântica do texto \cite{toraman}.
    \begin{code}
    \captionof{listing}{Função de Tokenização}
    \label{f:tokenization}
    \begin{mintedbox}{python}
    def tokenize_dataset(model_name: str):
        t = transformers.AutoTokenizer.from_pretrained(model_name, use_fast=True)
        t.pad_token = t.eos_token
        data_collator = transformers.DataCollatorForLanguageModeling(t, mlm=False)
    
        return data_collator, t
    \end{mintedbox}
    \end{code}
\end{itemize}

O fluxo de aplicação de pré-processamento segue a seguinte linha: transformação do conjunto de caracteres da entrada em letra minúscula, remoção de espaços redundantes e finalmente a tokenização dos caracteres. O Trecho de Código \ref{f:prepare} encapsula as etapas de transformação e remoção de espaços, sendo seguida pela aplicação da função de tokenização.

\begin{code}
\captionof{listing}{Aplicação das Etapas de Pré-processamento}
\label{f:prepare}
\begin{mintedbox}{python}
def prepare_dataset(df: pd.DataFrame) -> pd.DataFrame:
    df['text'] = df['text'].apply(text_to_lower)
    df['ethical_us'] = df['ethical_us'].apply(text_to_lower)
    df['text'] = df['text'].apply(pad_multiple_spaces)
    df['ethical_us'] = df['ethical_us'].apply(pad_multiple_spaces)
    instruction = construct_instruction()
    df['data'] = df.apply((lambda row: apply_in_row(row, instruction)), axis=1)
    
    return df

df = prepare_dataset(df)
data_collator, t = tokenize_dataset(model_name)
\end{mintedbox}
\end{code}

% split de treino e teste
A construção de um modelo de \textit{deep learning} requer a utilização de uma quantidade considerável de dados. No entanto, a utilização de dados que não tenham sido corretamente separados pode resultar em um desafio notável no treinamento, conhecido como \textit{overfitting}. Para Valdenegro-Toro \& Sabatelli \cite{overfitting}, overfitting pode ser definido como a falta de generalização em um modelo de \textit{machine learning}.

Portanto, após o pré-processamento, os dados foram divididos em dois conjuntos separados: um conjunto de treino e um conjunto de teste. Isto foi realizado através da utilização da técnica de validação \textit{hold-out}. Yadav \& Shuckla \cite{yadav} definem a técnica \textit{hold-out} como:
\blockquote{\textit{Hold-out validation} foi proposta para eliminar o problema de \textit{overfitting} que existia na validação de re-substituição. Os dados são divididos em duas partes que não se sobrepõem e estas duas partes são utilizadas para treinar e testar, respectivamente. A parte que é utilizada para o teste é a parte de \textit{hold-out}. Tem este nome porque retiramos essa parte para teste e  o modelo aprende utilizando a parte restante dos dados \cite{yadav}.}

A proporção de dados atribuídos a teste e treino pode variar. Normalmente, é utilizada uma divisão de 20\% de teste e 80\% de treino, que pode ser ajustada de acordo com o modelo específico. No caso de uma porcentagem de 10\%, existe um risco de \textit{overfitting} \cite{yadav}. Portanto, foi selecionada uma porcentagem de 20\% para os dados de teste neste estudo. No Trecho de Código \ref{f:div} é observada esta divisão dentro do código do modelo.

\begin{code}
\captionof{listing}{Divisão de Dados de Treino e Teste}
\label{f:div}
\begin{mintedbox}{python}
def train_test_split(df: pd.DataFrame) -> pd.DataFrame:
    df = datasets.Dataset.from_pandas(df)
    df = df.train_test_split(test_size=0.2)

    return df
\end{mintedbox}
\end{code}

\section{Desenvolvimento do Modelo}

O processo de desenvolvimento do modelo foi composto de um conjunto de etapas que, em ordem, contribuíram para a definição do modelo utilizado, treinamento por \textit{fine-tuning} e hiperparametrização. Em cada etapa, os resultados foram submetidos a uma avaliação, o que permitiu a aplicação de modificações em caso de necessário e em tempo real. A primeira etapa consistiu na escolha do modelo utilizado, e para isso, foi utilizada a plataforma Hugging Face\footnote{\url{https://huggingface.co/models}}, que possui um grande catálogo de modelos pré-treinados para diversas finalidades, incluindo PLN. 

A fim de identificar o modelo mais adequado, foi estabelecida uma série de modelos para efeitos de avaliação, utilizando o dataset deste projeto. Este conjunto inclui os seguintes modelos: Falcon40B \cite{falcon40b}, Falcon7B \cite{falcon40b}, BERT \cite{bert}, RoBERta \cite{roberta} e Mistral7B \cite{mistral}. Os modelos foram avaliados de acordo com os parâmetros definidos na Tabela \ref{tab:params} e o modelo Mistral7B \cite{mistral} demonstrou o melhor desempenho, produzindo as respostas mais otimizadas sobre o tema. Cada parâmetro possui um conjunto de valores associado que foi utilizado para alcançar os resultados. O processo de treinamento por \textit{fine-tuning} dos modelos foi realizado na plataforma Google Collab\footnote{\url{https://colab.research.google.com/}}. Os parâmetros e os valores foram definidos por serem considerados os básicos e padrões de um modelo de \acrshort{IA} \cite{hyper1, hyper2}. Os valores para \textit{batch size} e \textit{number of epochs} foram ajustados para serem inferiores aos valores padrão devido aos recursos computacionais disponíveis para treinar os modelos. No Trecho de Código \ref{f:model} podem ser observadas as configurações do processo de \textit{fine-tuning} do modelo.

\begin{table}[ht]
    \centering
    \caption{Parâmetros definidos para o teste dos modelos}
    \begin{tabular}{|l|l|l|}
    \hline
         \textbf{Parâmetro} & \textbf{Valores Testados} & \textbf{Valores Finais} \\ \hline
         \textit{Learning Rate} & 1e-2, 1e-3, 1e-4 e 1e-5 & \textbf{1e-3} \\ \hline
         \textit{Number of Epochs} & 3, 4, 5 e 6 & \textbf{3} \\ \hline
         \textit{Batch Size} & 4, 5 e 6 & \textbf{4} \\ \hline
         \textit{Weight Decay} & 1e-1, 1e-2 e 1e-3 & \textbf{1e-1} \\ \hline
    \end{tabular}
    \label{tab:params}
\end{table}

\begin{code}
\captionof{listing}{Aplicação da Técnica QLORA}
    \label{f:model}
    \begin{mintedbox}{python}
    def train_model(model, lr, batch_size, num_epochs, tokenized_data, collator):
        ... # QLORA steps
        
        training_args = transformers.TrainingArguments(
            output_dir= "checkpoints_output",
            learning_rate=lr,
            per_device_train_batch_size=batch_size,
            per_device_eval_batch_size=batch_size,
            num_train_epochs=num_epochs,
            weight_decay=0.01,
            logging_strategy="epoch",
            evaluation_strategy="epoch",
            save_strategy="epoch",
            load_best_model_at_end=True,
            warmup_steps=2,
            fp16=True,
            optim="paged_adamw_8bit",
        )
    
        trainer = transformers.Trainer(
            model=model,
            train_dataset=tokenized_data["train"],
            eval_dataset=tokenized_data["test"],
            args=training_args,
            data_collator=collator
        )
    
        # train model
        model.config.use_cache = False  # silence the warnings
        trainer.train()
    
        # renable warnings
        model.config.use_cache = True
    
        return model
\end{mintedbox}
\end{code}

% QLORA
Como mencionado nos parágrafos anteriores, o método de \textit{fine-tuning} foi utilizado para o treinamento do modelo. No entanto, é necessário ter em consideração que o treino é um processo custoso, uma vez que muitos modelos requerem mais de 30 GB de memória GPU. Uma das metodologias mais comuns para contornar este problema é a quantização. No entanto, estas técnicas só são eficazes durante a fase de inferência e são pouco úteis durante a fase de treino \cite{qlora}. 

Para preencher esta lacuna foi proposto o método QLORA, que utiliza uma técnica para quantizar um modelo pré-treinado para quatro bits e, em seguida, adiciona um pequeno conjunto de pesos adaptáveis \cite{qlora}. Estes pesos são ajustados por gradientes durante a \emph{backpropagation} através dos pesos quantizados \cite{qlora}.
O QLORA introduz uma combinação de três componentes. O primeiro é a quantização utilizando NormalFloat de quatro bits, que é teoricamente o tipo ótimo para dados normalmente distribuídos \cite{qlora}. O segundo é a quantização dupla, que funciona ao quantificar as constantes de quantização, economizando uma média de cerca de 0,37 bits por parâmetro \cite{qlora}. Por fim, são introduzidos os otimizadores de paginação, que usam a memória unificada NVIDIA para evitar os picos de memória de pontos de controle de gradiente ocorridos ao processar um \textit{mini-batch} de sequência longa \cite{qlora}. O Trecho de Código \ref{f:qlora} ilustra a forma como o método é utilizado no código do modelo.

\begin{code}
\captionof{listing}{Aplicação da Técnica QLORA}
    \label{f:qlora}
    \begin{mintedbox}{python}
    def train_model(model, lr, batch_size, num_epochs, tokenized_data, collator):
        model.train() # training state
        model.gradient_checkpointing_enable()
        model = peft.prepare_model_for_kbit_training(model) # turn into qlora
    
        # lora config
        config = peft.LoraConfig(
            r=32,
            lora_alpha=64,
            target_modules=["q_proj"],
            lora_dropout=.1,
            bias="none",
            task_type="CAUSAL_LM"
        )
        config.inference_mode = False
    
        model = peft.get_peft_model(model, config) # model in lora style
    
        ... # training steps
    \end{mintedbox}
\end{code}

\section{\textit{Retrieval-Augumented Generation}}

\textcolor{red}{A ser desenvolvido após a qualificação} 

\section{Validação da Ferramenta Proposta} 

\textcolor{red}{A ser desenvolvido após a qualificação} 

\section{Discussões}
 
As conclusões iniciais deste estudo indicam uma falta de ferramentas adequadas para a aplicação da ética em \acrshort{IA} em sistemas reais. Esta conclusão é exposta no Capítulo \ref{ref}, que apresenta em detalhe os resultados da revisão sistemática da literatura realizada. De forma a colaborar com aplicação da ética em \acrshort{IA}, foi proposta uma ferramenta baseada em LLM que realiza a tradução de requisitos éticos em historias de usuário éticas.
Foram identificadas diversas ferramentas que utilizam LLMs no contexto de historias de usuário. Zhang et al. \cite{zheying} explora o desenvolvimento de uma ferramenta que automatize a melhoria da qualidade de histórias de usuário através da integração de agentes baseados em LLM em ambientes reais. Luitel et al. \cite{luitel} desenvolveram uma maneira de utilizar um sistema baseado em \emph{Large Language Models} (LLMs) para gerar previsões contextualizadas para o preenchimento de requisitos de software. Estes estudos serviram como base para a ferramenta.

Neste contexto, a ferramenta proposta serve como um meio prático de ligação entre a ética em \acrshort{IA} e a engenharia de requisitos, permitindo a automatização deste processo. Halme et al. \cite{halme} propuseram a noção da utilização de histórias de usuário éticas como meio de facilitar a integração da ética em \acrshort{IA} nas práticas convencionais de engenharia de software. E com base nessa proposta o modelo foi desenvolvido. 

Durante o desenvolvimento da ferramenta, ficou evidente que a disponibilidade limitada de dados relacionados a ética em \acrshort{IA} era um desafio para o treinamento do modelo. Para resolver este problema, foi utilizado o processo de \textit{data augmentation}, se revelou necessário para permitir que o modelo fosse treinado dentro do prazo limitado deste trabalho. O modelo treinado demonstrou um desempenho satisfatório, mas para isso foi necessário que as predições fossem limitadas apenas a descrições de histórias de usuário. Portanto o modelo ainda não as gera no \textit{template} completo proposto por Halme et al. \cite{halme}. O formato dos resultados gerados até o momento é:
\begin{center}    
"como <persona> quero <fazer algo> <para que>".
\end{center}
